\section{Introduction}
\label{sec:intro}
The International Roughness Index (IRI) is the standard summary of longitudinal road roughness~\cite{sayers1986tp45,astm_e1926}.
Measuring it requires a longitudinal profile, typically from an inertial laser profilometer, which is too expensive to survey minor roads frequently.
Ordinary vehicles that already carry accelerometers---smartphones, telematics units, connected cars---could close this gap if their vibration could be mapped to IRI~\cite{sattar2018review,sang2024iri,levenberg2021lira}.
Two obstacles recur: labelled data are scarce because each label requires a profilometer run, and the mapping from vibration to IRI depends on vehicle, mounting and driving speed~\cite{sayers1986tp46,thilak2021,mounting2026}.
Simulation is attractive for the first obstacle because the label of a simulated road is known exactly.

This paper started as a revision of our own earlier, unpublished study, which trained a CNN on vibration from the BeamNG.tech game engine, reported an MAE of 1.65\,m/km on a simulated test track, and claimed speed invariance and real-world validity from a crowdsourced dataset without IRI labels.
Re-examining that work revealed problems that are easy to make and hard to see in the final numbers: the labels were not IRI, the test track was also used for early stopping and calibration, test predictions above a threshold were discarded, and the crowdsourced labels were explained by driving speed.
Rather than polish that result, we asked which conclusions survive a protocol that excludes these failure modes, and we moved the main evaluation to public car data with profilometer labels.

We make four contributions.
(i) An IRI implementation that reproduces the P79 profilometer's reported IRI on public data ($r=$\,\iriValR{} over \iriValN{} segments), and evidence that the heuristic label correction used for game-engine roads lowers the IRI of real urban roads by \legacyCorrUrbanMape\,\% (Section~\ref{sec:res_labels}).
(ii) A quantification of how test-domain calibration and prediction filtering shaped our earlier results (Section~\ref{sec:res_protocol}).
(iii) A route-disjoint comparison of constant, speed-only, vibration-index, gradient-boosting, MLP and CNN estimators on \liraKm\,km of real car data with profilometer IRI, with segment-level uncertainty, and an analysis of speed as a confounder (Sections~\ref{sec:res_real}--\ref{sec:res_speed}).
(iv) Zero-shot and few-shot transfer from a randomised half-car and sensor simulator to real cars, and controlled tests of speed invariance, unseen vehicles and sensor robustness with criteria fixed in advance (Sections~\ref{sec:res_sim}--\ref{sec:res_controlled}).

In brief, simulation over measured profiles nearly matches real training (\zsPsevenGbmMae{} vs.\ \liraGbmMae\,m/km), but no accurate estimator meets our unseen-vehicle criterion, fixed in advance.
Code, protocol, configuration and all predictions are provided with the paper.
